\begin{frame}{Lecture summary}
  \small
  \begin{itemize}\itemsep0.22em
    \item[\textcolor{red}{$\blacktriangleright$}] \textbf{Motivation:} video is the natural domain for dynamics and ``world'' hypotheses; we contrast \textbf{pixel generative} stacks with \textbf{latent predictive} models.
    \item[\textcolor{red}{$\blacktriangleright$}] \textbf{World Models:} Can be used to generate interactive videos, 3D scenes and latent world representations.
    \item[\textcolor{red}{$\blacktriangleright$}] \textbf{The Simulation Gap:} Despite stunning visuals (e.g., pouring liquid), current models lack symbolic or true physics-based understanding (e.g., accurate fluid dynamics).
    \item[\textcolor{red}{$\blacktriangleright$}] \textbf{Takeaway:} Mathematical rigor (JEPA loss, DiT diffusion), hands-on data collection (GameNGen, Genie), and grounding in both visual and physics-based techniques will empower the next generation of world-model researchers and practitioners.
    \item[\textcolor{red}{$\blacktriangleright$}] \textbf{Looking ahead:} Mastery of foundational world models is key to enabling autonomous systems and wearable AI assistants as the field advances from ``one-minute clips'' to unlimited, open-ended environments—and toward the frontier of Artificial General Intelligence.    
  \end{itemize}
\end{frame}
